\documentclass[fleqn]{beamer}

\mode<presentation> {
\usetheme{Madrid}
\usecolortheme{dolphin}
}

\usepackage[utf8]{inputenc}
\usepackage[brazilian]{babel}
\usepackage{amsmath}
\usepackage{amssymb}
\usepackage{booktabs}
\usepackage{listings}
\usepackage{xcolor}
\usepackage{tikz}

\lstset{
	basicstyle=\ttfamily\scriptsize,
	breaklines=true,
	keywordstyle=\color{blue!60!black},
	commentstyle=\color{gray},
	stringstyle=\color{green!45!black},
	showstringspaces=false,
	columns=fullflexible,
	frame=single,
	framesep=4pt,
	aboveskip=6pt,
	belowskip=2pt,
}

\title[Aula 5 -- Caixeiro Viajante]{O Problema do Caixeiro Viajante}
\subtitle{Formulações MTZ e de Fluxo Multi-Produto}
\author[Santi]{Prof. \'Everton Santi}
\institute[UFRN/ECT]{ECT2524 -- Introdução à Otimização \\ Escola de Ciências e Tecnologia -- UFRN}
\date{Aula 5}

\begin{document}

\begin{frame}
	\titlepage
\end{frame}

\begin{frame}
	\frametitle{Sumário}
	\tableofcontents
\end{frame}

%------------------------------------------------------------
\section{Contexto}
%------------------------------------------------------------

\begin{frame}
	\frametitle{O Problema do Caixeiro Viajante}
	\begin{itemize}
		\item \textbf{Travelling Salesman Problem} (TSP): um dos problemas
			mais estudados em otimização combinatória;
		\item Dado um conjunto de cidades e a distância entre cada par,
			encontrar a rota que:
		\begin{itemize}
			\item parte de uma cidade;
			\item visita todas as demais exatamente uma vez;
			\item retorna à cidade de origem;
			\item minimiza a distância total percorrida;
		\end{itemize}
		\item Fácil de \textbf{enunciar}, difícil de \textbf{resolver}: o
			número de rotas possíveis cresce como $(n-1)!/2$.
	\end{itemize}
\end{frame}

\begin{frame}
	\frametitle{Uma rota}
	\begin{center}
	\begin{tikzpicture}[scale=0.95,
		cid/.style={circle,draw,fill=blue!15,minimum size=6mm,inner sep=0pt}]
		\node[cid] (c1) at (0.5,2.2) {1};
		\node[cid] (c2) at (2.6,3.6) {2};
		\node[cid] (c3) at (5.2,3.2) {3};
		\node[cid] (c4) at (6.6,1.2) {4};
		\node[cid] (c5) at (4.4,0.2) {5};
		\node[cid] (c6) at (1.9,0.6) {6};
		\draw[blue!60,thick,->] (c1) -- (c2);
		\draw[blue!60,thick,->] (c2) -- (c3);
		\draw[blue!60,thick,->] (c3) -- (c4);
		\draw[blue!60,thick,->] (c4) -- (c5);
		\draw[blue!60,thick,->] (c5) -- (c6);
		\draw[blue!60,thick,->] (c6) -- (c1);
	\end{tikzpicture}
	\end{center}
	\begin{itemize}
		\item Uma rota válida visita cada cidade uma vez e fecha o ciclo;
		\item O objetivo compara o \textbf{comprimento total} de rotas
			como esta.
	\end{itemize}
\end{frame}

%------------------------------------------------------------
\section{Duas formulações}
%------------------------------------------------------------

\begin{frame}
	\frametitle{Duas formas de eliminar subrotas}
	\begin{itemize}
		\item Toda formulação de PLI para o TSP precisa, além de escolher
			os arcos da rota, impedir \textbf{subrotas} -- vários ciclos
			desconexos cobrindo só parte das cidades;
		\item Hoje vemos \textbf{duas} formulações que fazem isso de
			jeitos diferentes:
		\begin{itemize}
			\item \textbf{MTZ} (Miller--Tucker--Zemlin) -- compacta, fácil
				de implementar, limitante inferior \textbf{fraco};
			\item \textbf{MCF} (Fluxo de Mercadorias Multi-Produto) --
				compacta, limitante inferior \textbf{forte}, mas com muito
				mais variáveis;
		\end{itemize}
		\item Ambas são \textbf{fornecidas} -- a tarefa é
			\textbf{implementá-las e comparar}, não derivá-las.
	\end{itemize}
\end{frame}

\begin{frame}
	\frametitle{Miller--Tucker--Zemlin (1960)}
	\begin{itemize}
		\item Uma das formulações de programação linear inteira mais
			conhecidas para o TSP;
		\item Vantagem: número \textbf{polinomial} de restrições e fácil
			de implementar;
		\item Limitação: o limitante inferior que ela produz é
			relativamente \textbf{fraco} -- é o que a MCF, vista mais
			adiante, corrige, ao custo de muito mais variáveis.
	\end{itemize}
\end{frame}

\begin{frame}
	\frametitle{Parâmetros e variáveis (MTZ)}
	Parâmetros:
	\begin{itemize}
		\item $n$ -- número de cidades;
		\item $d_{ij}$ -- distância entre as cidades $i$ e $j$, $\forall
			i,j = 1,\ldots,n$;
	\end{itemize}
	Variáveis de decisão:
	\begin{itemize}
		\item $x_{ij} \in \{0,1\}$ -- vale 1 se a rota vai da cidade $i$
			para a cidade $j$, $\forall i,j = 1,\ldots,n$;
		\item $u_i \in \{1,\ldots,n-1\}$ -- posição da cidade $i$ na rota,
			$\forall i = 2,\ldots,n$ (a cidade 1 é sempre a origem/retorno,
			por isso não precisa de $u_1$).
	\end{itemize}
\end{frame}

\begin{frame}
	\frametitle{Função objetivo e restrições de grau}
	$$
		\text{minimize}~Z = \sum_{i=1}^{n}\sum_{j=1,j\neq i}^{n} d_{ij}x_{ij}
	$$
	sujeito a:
	$$
		\sum_{i=1,i\neq j}^{n} x_{ij} = 1,~\forall~j = 1,\ldots,n
	$$
	$$
		\sum_{j=1,j\neq i}^{n} x_{ij} = 1,~\forall~i = 1,\ldots,n
	$$
	\begin{itemize}
		\item A primeira igualdade garante que toda cidade $j$ é
			\textbf{alcançada} a partir de exatamente uma outra;
		\item A segunda garante que toda cidade $i$ tem exatamente uma
			\textbf{saída};
		\item Sozinhas, porém, essas restrições não impedem uma solução
			partida em vários pedaços.
	\end{itemize}
\end{frame}

\begin{frame}
	\frametitle{O problema das subrotas}
	Suponha $n=4$ e a solução:
	$$
		x = \begin{pmatrix}
			0 & 0 & 1 & 0 \\
			0 & 0 & 0 & 1 \\
			1 & 0 & 0 & 0 \\
			0 & 1 & 0 & 0
		\end{pmatrix}
	$$
	\begin{itemize}
		\item Toda linha e toda coluna somam 1 -- as restrições de grau
			estão satisfeitas;
		\item Mas a solução descreve \textbf{duas} rotas separadas: $1
			\to 3 \to 1$ e $2 \to 4 \to 2$, cada uma cobrindo só parte das
			cidades;
		\item Precisamos de restrições que \textbf{proíbam} esse tipo de
			solução, deixando só o ciclo único com as $n$ cidades.
	\end{itemize}
\end{frame}

\begin{frame}
	\frametitle{Duas subrotas \emph{vs.} uma rota}
	\begin{columns}
	\begin{column}{0.48\textwidth}
	\centering
	\begin{tikzpicture}[scale=0.7,
		cid/.style={circle,draw,fill=red!12,minimum size=6mm,inner sep=0pt}]
		\node[cid] (c1) at (0,1.6) {1};
		\node[cid] (c3) at (1.8,1.6) {3};
		\node[cid] (c2) at (0,0) {2};
		\node[cid] (c4) at (1.8,0) {4};
		\draw[red!70!black,thick,->] (c1) to[bend left=25] (c3);
		\draw[red!70!black,thick,->] (c3) to[bend left=25] (c1);
		\draw[red!70!black,thick,->] (c2) to[bend left=25] (c4);
		\draw[red!70!black,thick,->] (c4) to[bend left=25] (c2);
	\end{tikzpicture}
	\\[0.3cm] \footnotesize inválida: duas subrotas
	\end{column}
	\begin{column}{0.48\textwidth}
	\centering
	\begin{tikzpicture}[scale=0.7,
		cid/.style={circle,draw,fill=blue!12,minimum size=6mm,inner sep=0pt}]
		\node[cid] (c1) at (0,1.6) {1};
		\node[cid] (c3) at (1.8,1.6) {3};
		\node[cid] (c2) at (0,0) {2};
		\node[cid] (c4) at (1.8,0) {4};
		\draw[blue!60,thick,->] (c1) -- (c3);
		\draw[blue!60,thick,->] (c3) -- (c4);
		\draw[blue!60,thick,->] (c4) -- (c2);
		\draw[blue!60,thick,->] (c2) -- (c1);
	\end{tikzpicture}
	\\[0.3cm] \footnotesize válida: um único ciclo
	\end{column}
	\end{columns}
	\vspace{0.3cm}
	Ambas satisfazem as restrições de grau -- só as restrições de
	eliminação de subrota descartam a da esquerda.
\end{frame}

\begin{frame}
	\frametitle{Restrições MTZ de eliminação de subrota}
	$$
		u_i - u_j + (n-1)x_{ij} \leq n-2,~\forall~i,j = 2,\ldots,n;~i\neq j
	$$
	$$
		1 \leq u_i \leq n-1,~\forall~i = 2,\ldots,n
	$$
	\begin{itemize}
		\item $u_i$ marca em que \textbf{posição} da rota a cidade $i$ é
			visitada -- se $u_3 = 4$, a cidade 3 é a quarta a ser visitada;
		\item Se $x_{ij}=1$ (a rota vai de $i$ para $j$), a restrição
			força $u_j \geq u_i + 1$: $j$ só pode vir \textbf{depois} de
			$i$ na sequência;
		\item Isso é impossível de satisfazer numa subrota que não passa
			pela cidade 1 -- daí a eliminação.
	\end{itemize}
\end{frame}

\begin{frame}
	\frametitle{Formulação MTZ completa}
	\begin{align*}
		\text{minimize}~Z &= \sum_{i=1}^{n}\sum_{j=1,j\neq i}^{n} d_{ij}x_{ij} \\
		\text{sujeito a: } \sum_{i=1,i\neq j}^{n} x_{ij} &= 1,~\forall j = 1,\ldots,n \\
		\sum_{j=1,j\neq i}^{n} x_{ij} &= 1,~\forall i = 1,\ldots,n \\
		u_i - u_j + (n-1)x_{ij} &\leq n-2,~\forall i,j = 2,\ldots,n;~i\neq j \\
		1 \leq u_i &\leq n-1,~\forall i = 2,\ldots,n \\
		x_{ij} &\in \{0,1\},~\forall i,j = 1,\ldots,n \\
		u_i &\in \mathbb{Z},~\forall i = 2,\ldots,n
	\end{align*}
\end{frame}

\begin{frame}
	\frametitle{Fluxo de Mercadorias Multi-Produto (MCF)}
	\begin{itemize}
		\item Ideia: a cidade 1 é uma central que envia \textbf{um produto
			diferente} para cada uma das outras $n-1$ cidades, usando só
			os arcos que a rota percorre;
		\item Se a rota se quebrasse em subrotas, o produto destinado a
			uma cidade de uma subrota sem a cidade 1 \textbf{nunca
			chegaria} -- daí a eliminação;
		\item Variável nova: $f_{ij}^{k} \geq 0$ -- quantidade do produto
			$k$ no arco $(i,j)$, $\forall i,j = 1,\ldots,n,~i\neq j$ e
			$\forall k = 2,\ldots,n$ (um produto por cidade destino).
	\end{itemize}
\end{frame}

\begin{frame}
	\frametitle{Conservação de fluxo: origem e destino}
	Para cada produto $k = 2,\ldots,n$:
	\begin{itemize}
		\item \textbf{Origem} -- a cidade 1 emite 1 unidade do produto
			$k$ (saída menos entrada $=1$):
	\end{itemize}
	$$
		\sum_{j=1,j\neq 1}^{n} f_{1j}^{k} - \sum_{j=1,j\neq 1}^{n} f_{j1}^{k} = 1,
		~\forall~k = 2,\ldots,n
	$$
	\begin{itemize}
		\item \textbf{Destino} -- a cidade $k$ absorve essa unidade
			(saída menos entrada $=-1$):
	\end{itemize}
	$$
		\sum_{j=1,j\neq k}^{n} f_{kj}^{k} - \sum_{j=1,j\neq k}^{n} f_{jk}^{k} = -1,
		~\forall~k = 2,\ldots,n
	$$
\end{frame}

\begin{frame}
	\frametitle{Conservação de fluxo: demais cidades, e acoplamento}
	\begin{itemize}
		\item \textbf{Demais cidades} -- toda $i \neq 1,k$ só repassa o
			que recebe (saída menos entrada $=0$):
	\end{itemize}
	$$
		\sum_{j=1,j\neq i}^{n} f_{ij}^{k} - \sum_{j=1,j\neq i}^{n} f_{ji}^{k} = 0,
		~\forall~i = 2,\ldots,n;~i \neq k;~\forall~k = 2,\ldots,n
	$$
	\begin{itemize}
		\item \textbf{Acoplamento com a rota} -- só há fluxo num arco que
			a rota realmente usa:
	\end{itemize}
	$$
		f_{ij}^{k} \leq x_{ij},~\forall i,j = 1,\ldots,n;~i\neq j
	$$
	\begin{itemize}
		\item Se $x_{ij}=0$, o acoplamento zera $f_{ij}^{k}$ para todo
			$k$ -- só a rota escolhida pode carregar fluxo.
	\end{itemize}
\end{frame}

\begin{frame}
	\frametitle{MTZ \emph{vs.} MCF}
	\begin{center}
	\footnotesize
	\begin{tabular}{lcc}
		\toprule
		 & MTZ & MCF \\
		\midrule
		Variáveis extras & $u_i$ ($O(n)$) & $f_{ij}^{k}$ ($O(n^3)$) \\
		Restrições extras & $O(n^2)$ & $O(n^3)$ \\
		Limitante inferior (relaxação) & fraco & forte \\
		Equivalente, em força, a & -- & Dantzig--Fulkerson--Johnson \\
		\bottomrule
	\end{tabular}
	\end{center}
	\begin{itemize}
		\item A MCF dá um limitante tão bom quanto o da formulação
			clássica com restrições exponenciais de subrota -- \textbf{sem}
			precisar de plano de cortes;
		\item O preço é o crescimento cúbico do número de variáveis --
			inviável para instâncias grandes.
	\end{itemize}
\end{frame}

%------------------------------------------------------------
\section{A tarefa (avaliativa)}
%------------------------------------------------------------

\begin{frame}
	\frametitle{Caráter avaliativo}
	\begin{itemize}
		\item A tarefa desta aula \textbf{vale nota}: peso de
			\textbf{25\%} na Unidade 1 da disciplina;
		\item As formulações MTZ e MCF (slides anteriores) são
			\textbf{fornecidas prontas} -- não é preciso derivá-las;
		\item O que se avalia é a \textbf{implementação} das duas, a
			execução nas instâncias corretas e a qualidade da comparação
			entre elas;
		\item Detalhes de formato de entrega e prazo: ver o enunciado
			completo na página da aula, no site do curso.
	\end{itemize}
\end{frame}

\begin{frame}
	\frametitle{O que é pedido}
	\begin{enumerate}
		\item Escrever as duas formulações em AMPL:
			\texttt{tsp\_mtz.mod} e \texttt{tsp\_mcf.mod};
		\item Escrever, em Python, a leitura de uma instância TSPLIB
			(coordenadas $\to$ matriz de distâncias $d_{ij}$);
		\item \textbf{Tentar} resolver as \textbf{duas formulações} nas
			\textbf{quatro} instâncias (8 combinações), com 1 thread e
			tempo limite de 10 min no solver;
		\item Reportar LB, UB, GAP, tempo e status numa única tabela (8
			linhas) -- registrando como \textbf{impraticável} qualquer
			combinação que travar ou estourar a memória.
	\end{enumerate}
\end{frame}

\begin{frame}
	\frametitle{As instâncias}
	\begin{center}
	\footnotesize
	\begin{tabular}{lc}
		\toprule
		Instância (TSPLIB) & Cidades \\
		\midrule
		Djibouti (\texttt{dj38}) & 38 \\
		Uruguai (\texttt{uy734}) & 734 \\
		Luxemburgo (\texttt{lu980}) & 980 \\
		Omã (\texttt{mu1979}) & 1\,979 \\
		\bottomrule
	\end{tabular}
	\end{center}
	\begin{itemize}
		\item Quatro instâncias reais da coleção \emph{National TSP}
			(formato TSPLIB, \texttt{EUC\_2D}), em ordem crescente de
			tamanho;
		\item As duas formulações são tentadas nas quatro -- não há
			instância reservada só para uma delas;
		\item Fonte e formato dos arquivos: ver o enunciado da tarefa, no
			site do curso.
	\end{itemize}
\end{frame}

\begin{frame}
	\frametitle{Quando a MCF fica impraticável}
	\begin{itemize}
		\item A MCF tem $O(n^3)$ variáveis de fluxo -- em Djibouti
			($n=38$) são dezenas de milhares; em Luxemburgo ($n=980$),
			já passa de \textbf{940 milhões}; em Omã ($n=1979$), de
			\textbf{7,7 bilhões};
		\item O limite de \textbf{10 min} vale para o \textbf{solver} --
			não protege a fase de \textbf{montar} o modelo, que pode
			travar ou esgotar a memória bem antes disso nas instâncias
			maiores;
		\item Isso é esperado: se acontecer, interrompam a execução,
			anotem o sintoma e registrem a combinação como
			\textbf{impraticável} na tabela -- não precisam forçar até o
			fim;
		\item Combinado ao outro lado do compromisso -- o GAP da MTZ
			\textbf{crescendo} com $n$, dentro do mesmo tempo -- a tabela
			completa mostra os dois limites na prática: a MCF fecha
			melhor, mas não escala; a MTZ escala, mas fecha pior.
	\end{itemize}
\end{frame}

%------------------------------------------------------------
\section{Referências}
%------------------------------------------------------------

\begin{frame}
	\frametitle{Referências}
	\begin{itemize}
		\item Miller, C. E.; Tucker, A. W.; Zemlin, R. A. \emph{Integer
			Programming Formulation of Traveling Salesman Problems}.
			Journal of the ACM, 1960.
		\item Wong, R. T. \emph{Integer Programming Formulations of the
			Travelling Salesman Problem}. IEEE ICCC, 1980 (formulação de
			fluxo multi-produto);
		\item TSPLIB / National TSP instances:
			\url{https://www.math.uwaterloo.ca/tsp/world/countries.html};
		\item AMPL: \url{https://ampl.com};
		\item Documentação do \texttt{amplpy}: \url{https://amplpy.ampl.com}.
	\end{itemize}
\end{frame}

\end{document}
